% =====================================================================
%  Results and Discussion — two case studies (IEEE 123 and IEEE 9500)
%  Drop-in section for an IEEE Transactions paper.
%  This file references (does NOT create) the already-inserted floats.
%
%  >>> Give each already-inserted float a UNIQUE \label:
%     IEEE 123 OpenDSS-validation table  ->  \label{tab:val123}
%     IEEE 9500 OpenDSS-validation table ->  \label{tab:val9500}
%     IEEE 123 DDDP performance table    ->  \label{tab:dddp123}
%     IEEE 9500 DDDP performance table   ->  \label{tab:dddp9500}
%     IEEE 123 gamma-sweep figure        ->  \label{fig:gamma123}
%     IEEE 9500 gamma-sweep figure       ->  \label{fig:gamma9500}
%     IEEE 123 DDDP linear LB figure     ->  \label{fig:conv123lin}
%     IEEE 123 DDDP ISOCP  LB figure     ->  \label{fig:conv123isocp}
%     IEEE 9500 DDDP linear LB figure    ->  \label{fig:conv9500lin}
%     IEEE 9500 DDDP ISOCP  LB figure    ->  \label{fig:conv9500isocp}
%     daily input-profiles figure        ->  \label{fig:profiles}
%  (your current files reuse \label{table3} and \label{tabledddp} twice,
%   which must be de-duplicated for the \ref{} calls below to resolve.)
% =====================================================================

\section{Results and Discussion}
\label{sec:results}

The framework is evaluated on the IEEE 123-bus and IEEE 9500-node unbalanced
three-phase distribution feeders. The results show that (i)~the lossless
LinDistFlow surrogate introduces systematic financial and physical errors that
only the ISOCP corrects, and (ii)~DDDP recovers the centralized ISOCP schedule
to within a fraction of a percent at every partition while reducing wall time by
up to $3.6\times$, making an otherwise $16.3$-minute solve operationally viable.

% ---------------------------------------------------------------
\subsection{Experimental Setup}
\label{ssec:setup}

The IEEE 9500-node feeder is reduced to a $2{,}752$-bus equivalent by referring
all secondary loads to the primary side of their distribution transformers.
On both feeders the top $50\%$ of load buses by nominal demand receive a
co-located PV unit and a two-hour battery, each sized at $45\%$ of the host
node's rated load. The planning horizon is $T\!=\!24$ hourly steps; battery
SoC continuity couples consecutive hours, requiring a full multi-period joint
optimization. Fig.~\ref{fig:profiles} shows the exogenous inputs. The
demand-linked price $c_t = 0.15\,\ell_t + 0.15~\text{\$/kWh}$ peaks in the
evening while PV peaks at noon, creating a charge/discharge arbitrage
opportunity that can only be fully exploited by coordinating the entire 24-hour
SoC trajectory.

Three OPF formulations are compared on a common network model: the lossless
LinDistFlow relaxation (\emph{Linear}), which drops the quadratic current terms
from the branch power balance and neglects ohmic losses; the
iteratively-tightened second-order cone program (\emph{ISOCP}), which adds
linearized cuts until the conic residual $P_{ij}^2 + Q_{ij}^2 - v_i\ell_{ij}$
falls below $10^{-4}$~p.u.\ on every branch; and, on the 123-bus feeder only,
the full \emph{Nonlinear} AC branch-flow model as an exact reference. All
convex programs use Gurobi~13/Pyomo/APPSI; every schedule is replayed
hour-by-hour through OpenDSS to verify physical fidelity.

% ---------------------------------------------------------------
\subsection{ISOCP Cut Parameter Calibration}
\label{ssec:gamma}

Each ISOCP cut over-satisfies the current cone violation by a factor
$(1{+}\gamma)$. Figs.~\ref{fig:gamma123} and~\ref{fig:gamma9500} sweep
$\gamma\!\in\![0.3,0.7]$; both curves are U-shaped with a minimum in
$[0.4,0.5]$. On the 123-bus feeder the initial cone gap of $1.47$ clears in
$8$ iterations at $\gamma\!\in\![0.4,0.5]$, against $12$--$14$ at the extremes.
The 9500-node feeder, with a larger initial gap of $5.9$ driven by the
heavily-loaded distribution transformer, sharpens the trend: $\gamma\!=\!0.3$
fails to converge within $31$ iterations while the minimum is at $\gamma\!=\!0.4$.
We fix $\gamma\!=\!0.5$ throughout---within the fast band on both feeders.

% =====================================================================
\subsection{IEEE 123-Bus Feeder}
\label{ssec:case123}

% ...............................................................
\subsubsection{OPF Formulation Comparison}
\label{sssec:123_valid}

Because LinDistFlow sets branch current magnitudes to zero in the power balance,
the substation need not supply any $I^2R$ loss in the model---but must on the
real feeder. This produces a systematic error in both the energy balance and the
cost objective that cannot be removed by post-processing.

Table~\ref{tab:val123} quantifies these effects. The per-bus DER dispatch
(active PV output and battery charge/discharge) is nearly identical between the
Linear and ISOCP solutions: individual decisions are governed by local prices and
capacity limits rather than network losses. The difference is systemic. The
Linear model plans for a 24-hour substation active import of $62{,}275$~kW while
OpenDSS requires $63{,}510$~kW---a $1{,}235$~kW shortfall equal to the entire
omitted loss, with a comparable reactive deficit. Worst-case per-phase, per-hour
deviations reach $1{,}234.9$~kW, $2{,}432.2$~kVAr, and $3.1\times10^{-3}$~p.u.
The Linear objective under-prices the schedule by \$342 ($2.0\%$): \$$17{,}101$
versus the AC-accurate \$$17{,}444$, the gap being exactly the cost of the
omitted loss.

The ISOCP corrects all of these. Its computed loss ($1{,}229.8$~kW) matches
OpenDSS ($1{,}230.9$~kW) to $0.1\%$; substation active import agrees to
$1.2$~kW; and worst-case per-phase deviations fall to $1.19$~kW, $42.1$~kVAr,
and $6\times10^{-4}$~p.u.---two to three orders of magnitude tighter. The ISOCP
converges in $9$ iterations and $18$~s and matches the Nonlinear AC reference
to \$0.14 ($17{,}443.55$ vs.\ $17{,}443.69$). Gurobi did not produce a
solution to the Nonlinear model after hours of runtime; IPOPT is used solely
as an accuracy check and produces a solution in approximately $40$~s.

% ...............................................................
\subsubsection{DDDP Performance}
\label{sssec:123_dddp}

Battery SoC continuity couples all $T\!=\!24$ steps into a single monolithic
problem; DDDP decomposes the horizon into $P$ disjoint windows and passes
future-cost information forward as Benders cuts. It maintains a monotone lower
bound $\underline{Z}^{(k)}\!=\!c_1(\bm{z}_1^{(k)})+\theta_1^{(k)}$ and an upper
bound $\overline{Z}^{(k)}\!=\!\sum_i c_i(\bm{z}_i^{(k)})$ that converge at
termination (Table~\ref{tab:dddp123},
Figs.~\ref{fig:conv123lin}--\ref{fig:conv123isocp}).

For both fidelities, $\underline{Z}^{(k)}$ converges monotonically to the same
optimum at every partition---\$$17{,}443.56$ (ISOCP) and \$$17{,}101.15$
(Linear)---with the optimality gap staying within a few thousandths of a percent
and never exceeding $0.13\%$ even at $P\!=\!12$. Partition count affects solve
time, not solution quality.

Increasing $P$ shrinks each window subproblem but lengthens the Benders chain,
since a cut propagates only one window per iteration (outer iterations:
$3, 6, 8, 11, 13, 16$ for $P\!=\!2$ to $12$, ISOCP). The sweet spot for the
ISOCP is $P\!=\!8$: three-step windows solve in $0.9$~s and total time falls
from $23.8$~s (centralized) to $11.4$~s ($2.1\times$). The ISOCP gains more than
the Linear model because full cone tightening is deferred to a single
post-convergence pass; only cheap inner cuts accompany each DDDP iteration. The
AC-accurate schedule of Table~\ref{tab:val123} is thus obtained for essentially
the cost of the relaxed decomposition. On this small feeder the Linear LP is
already $1.6$~s and overhead-dominated, limiting the improvement; the mechanism
pays in full on the 9500-node system.

% =====================================================================
\subsection{IEEE 9500-Node Feeder}
\label{ssec:case9500}

% ...............................................................
\subsubsection{OPF Formulation Comparison}
\label{sssec:9500_valid}

On the $2{,}752$-bus system neither Gurobi nor IPOPT produces a feasible
solution to the Nonlinear branch-flow model, making the ISOCP the only route
to an AC-accurate schedule.

Table~\ref{tab:val9500} shows this. The Linear model under-counts substation
active import by $1{,}960$~kW ($238{,}045$ vs.\ $240{,}005$~kW) and the
reactive shortfall exceeds $10{,}000$~kVAr. The Linear objective under-prices the schedule by \$499 (\$$65{,}358$
versus the ISOCP's \$$65{,}857$), the gap being the energy cost of the
$1{,}960$~kW of active loss the lossless model omits. This is larger than the \$342 gap on the 123-bus feeder, as network losses
increase with network size.
Worst-case per-phase, per-hour
deviations reach $1{,}959$~kW and $10{,}542$~kVAr, large enough to cause
systematic capacitor bank mis-operation, incorrect tap commands, and undervoltage
violations at distal buses. The ISOCP holds its accuracy at scale: substation
active and reactive power agree with OpenDSS to $3.1$~kW and $16.5$~kVAr;
losses match to $0.1\%$; and worst-case deviations are $3.08$~kW, $16.45$~kVAr,
and $1.3\times10^{-4}$~p.u.\ The reactive-power agreement is particularly
significant---reactive balance is the most sensitive indicator of unbalance
modelling accuracy---confirming the schedule is physically deliverable on a
heavily unbalanced feeder with nearly $10{,}000$ nodes.

% ...............................................................
\subsubsection{DDDP Performance}
\label{sssec:9500_dddp}

The centralized 24-hour ISOCP requires $977.6$~s ($16.3$~min), exceeding
day-ahead scheduling budgets. DDDP is therefore an operational necessity on this
feeder, not an optional refinement.

For both fidelities, $\underline{Z}^{(k)}$ tracks the centralized optimum
monotonically at every partition: the optimality gap is $0.006\%$ at $P\!=\!6$
and at most $0.055\%$ at the finest split (Table~\ref{tab:dddp9500},
Figs.~\ref{fig:conv9500lin}--\ref{fig:conv9500isocp}).

Unlike the 123-bus case, both fidelities benefit at every partition because each
window subproblem is large enough that shrinking it outweighs the additional
Benders iterations. The sweet spot is $P\!=\!6$: six four-step windows solve in
parallel in $23$~s each---one-seventh of the $163$~s for a centralized
single-window share---cutting ISOCP time from $977.6$~s to $272.3$~s
($3.6\times$). The Linear model falls from $50.7$~s to $16.6$~s ($3.0\times$)
over $P\!=\!6$--$8$; at this scale each LinDistFlow window is a large sparse LP
whose solve, not setup overhead, dominates. Beyond the sweet spot the lengthening
Benders chain ($4, 6, 8, 12, 16, 20$ iterations for $P\!=\!2$ to $12$) overtakes
the savings: time rises and the gap loosens to $0.03\%$ at $P\!=\!8$ and $0.05\%$
at $P\!=\!12$. As on the small feeder the ISOCP gains more, because full
tightening is deferred to one post-convergence sweep; the AC-accurate schedule is
obtained for essentially the cost of the relaxed parallel decomposition.

% ---------------------------------------------------------------
\subsection{Comparison with Temporal ADMM}
\label{ssec:tadmm}

We benchmark DDDP-OTD against the temporal ADMM (T-ADMM) of
Pinto~et~al.~\cite{PiBe20}, a representative decomposition for multi-period
OPF. T-ADMM partitions the horizon into $P$ windows and enforces inter-window
SoC continuity through an augmented Lagrangian, alternating between local
window solves and dual updates until the boundary consensus residuals fall
below tolerance. Table~\ref{tab:tadmm} compares the two methods on both
formulations and feeders at representative partitions.

\begin{table}[t]
\centering
\caption{Temporal ADMM~\cite{PiBe20} vs.\ DDDP-OTD.
         Obj.\ is the 24-hour schedule cost (\$/day).}
\label{tab:tadmm}
\small
\setlength{\tabcolsep}{3.5pt}
\begin{tabular}{l l c rrr rrr}
\toprule
 & & & \multicolumn{3}{c}{\textbf{Temporal ADMM}} & \multicolumn{3}{c}{\textbf{DDDP-OTD}} \\
\cmidrule(lr){4-6}\cmidrule(lr){7-9}
System & Model & $P$ & Obj.\ (\$/day) & Iters & Time (s) & Obj.\ (\$/day) & Iters & Time (s) \\
\midrule
\multirow{6}{*}{123-bus}
 & \multirow{3}{*}{Linear}
   & 4  & 17{,}101.0 & 35 &  4.1 & 17{,}101.2 &  5 & 1.1 \\
 & & 8  & 17{,}101.0 & 53 &  3.4 & 17{,}101.2 &  9 & 1.1 \\
 & & 12 & 17{,}100.8 & 73 &  3.8 & 17{,}102.4 & 13 & 1.2 \\
\cmidrule{2-9}
 & \multirow{3}{*}{ISOCP}
   & 4  & 17{,}442.8 & 33 & 16.0 & 17{,}443.7 &  8 & 12.1 \\
 & & 8  & 17{,}443.9 & 37 & 14.0 & 17{,}443.4 & 13 & 11.4 \\
 & & 12 & 17{,}444.1 & 44 & 15.3 & 17{,}465.5 & 16 & 12.1 \\
\midrule
\multirow{5}{*}{9500-node}
 & \multirow{2}{*}{Linear}
   & 2 & 65{,}356.7 &  59 & 279.9 & 65{,}357.6 &  3 & 29.8 \\
 & & 3 & 65{,}357.1 & 103 & 306.2 & 65{,}363.8 &  4 & 22.1 \\
\cmidrule{2-9}
 & \multirow{3}{*}{ISOCP}
   & 4  & 65{,}854.9 & 109 &  968.1 & 65{,}858.9 &  8 & 299.8 \\
 & & 8  & 65{,}852.9 & 114 &  734.8 & 65{,}875.7 & 16 & 288.9 \\
 & & 12 & 65{,}852.8 & 129 &  950.8 & 65{,}893.5 & 20 & 297.3 \\
\bottomrule
\end{tabular}
\end{table}

Both methods reproduce the centralized schedule cost to within $0.13\%$ at
every partition---T-ADMM to within $0.01\%$---so they are equivalent in accuracy
and differ only in speed. That difference is large, and it originates in
iteration count. T-ADMM takes $33$--$73$ iterations on the 123-bus feeder and
$59$--$129$ on the 9500-node feeder, whereas DDDP converges in $5$--$16$ and
$3$--$20$; both grow with $P$ as each window boundary adds a coupling
constraint, but T-ADMM stays several-fold higher throughout. The gap is
inherent to how each method enforces SoC continuity. T-ADMM satisfies the
boundary-continuity constraints only in the limit: it penalizes the boundary
mismatches and updates the dual prices once per iteration, and because ADMM
converges at a first-order (linear) rate, these primal and dual residuals decay
only geometrically toward the $10^{-4}$ tolerance---so many rounds are needed
before every boundary agrees. DDDP, by contrast, enforces continuity exactly on
each pass and instead refines a Benders approximation of the future cost-to-go;
each cut carries downstream cost information over the whole remaining horizon,
letting the windows anticipate their coupling directly and converge in far fewer
passes.

Because the iteration gap is fixed, wall time follows subproblem cost. On the
123-bus feeder, where each window solves in well under a second, T-ADMM is only
$3.1$--$3.7\times$ slower than DDDP for the Linear model and $1.2$--$1.3\times$
for the ISOCP. On the 9500-node feeder, where each window is a large ISOCP or
LP, the same iteration overhead becomes decisive: the Linear T-ADMM runs
$5.5$--$6.0\times$ slower than even the centralized solve---and $9$--$14\times$
slower than DDDP---while the ISOCP T-ADMM is $2.5$--$3.2\times$ slower than DDDP
at every partition. DDDP's advantage is thus largest precisely where it matters
most: on the large feeder, where the centralized solve is already too slow for
day-ahead use.


% ---------------------------------------------------------------
\subsection{Discussion}
\label{ssec:disc}

A key observation across both feeders is that the per-bus DER dispatch is nearly
identical between the Linear and ISOCP solutions---individual schedules are
governed by local prices and capacity limits rather than network losses. The
entire difference in cost and energy balance arises from whether $I^2R$ losses
are priced into the optimization, and that error grows with network size.

The ISOCP and DDDP together achieve what neither does alone. Deferring full cone
tightening to a single post-convergence pass amortizes its cost over the Benders
loop instead of repeating it at every outer iteration. The resulting schedule
matches the centralized ISOCP to within $0.006\%$ at $P\!=\!6$ on the 9500-node
feeder and is computed in $272.3$~s rather than $16.3$~min ($3.6\times$), with
monotone bounds and optimality gaps below $0.06\%$ at every tested partition on
both systems.

% ---------------------------------------------------------------
\section{Conclusion}
\label{sec:conclusion}

On unbalanced three-phase distribution feeders, LinDistFlow ignores network
losses and yields systematically mis-priced multi-period schedules, while
the full nonlinear branch-flow model is intractable at scale, leaving the
MPISOCP as the only AC-accurate tractable formulation.
This work presents a fully parallel temporal decomposition algorithm based on
DDDP for solving monolithic multi-period OPF problems under both linear and
MPISOCP formulations by partitioning the optimization horizon into
non-overlapping windows coordinated through inter-window information exchange.
A dual convergence criterion combining the relative optimality gap and a
lower-bound stability check ensured robust termination irrespective of the
number of partitions. Numerical results on the modified IEEE 123-bus and
IEEE 9500-node systems confirmed that the algorithm matches the centralized
global optimum across all partition counts, with an optimal partition size
balancing per-iteration cost against convergence speed. The parallel DDDP
strategy further reduced runtime substantially, especially for the non-convex
problems, by limiting full second-order cone solves to a single terminal
iteration after convergence.
